\section{The oriented projection over the completed field}
\label{app:oriented-projection}

We prove \cref{lem:completed-oriented-map}, including the continuous
identity that identifies the \(E_2\) tensor factor with connected base
change. The main construction is the split section in
\eqref{eq:oriented-section-diagram}. Here we verify its residue tags,
orientations, and family coherence on the parameter rings
\(R_T=C_{\cts}(T,F)\), where \(F=\widehat L(*)\) is the ordinary
completed field, and then pass to solid spectra.

\begin{proof}[Proof of \cref{lem:completed-oriented-map}]
The forgetful transformation in
\cite[Proposition~6.2.2 and Remark~6.2.3]{LurieEllipticII2018} goes from
deformations of a full group to deformations of the left term of a tagged
exact sequence with étale quotient. Corepresentability therefore gives a
map from the connected deformation ring to the full deformation ring.
We construct the reverse oriented map by a section of that transformation.

The field \(F\) is perfect of characteristic \(p\), and its inverse
Frobenius is continuous. Hence pointwise inverse Frobenius preserves
\(C_{\cts}(T,F)\), so \(R_T\) is perfect. For nonempty \(T\), the map
\(F\to R_T\) is flat, since \(F\) is a field, and relatively perfect, since
both Frobenius maps are isomorphisms. The ring \(R_T\) need not be a field
or Noetherian. The field \(F\) is Noetherian and \(F\)-finite.
More generally, for any perfect \(\F_p\)-algebra, every \(p\)-divisible group is nonstationary
and the absolute cotangent complex is almost perfect, by
\cite[Remark~3.4.2]{LurieEllipticII2018}. In particular the full group
over \(F\), without a connectedness assumption, is an eligible
Noetherian LKN witness for \((R_T,\Gamma_T)\), in the sense of
\cite[Definition~4.3.7 and Example~4.6.4]{BMR2026}. The connected pair has
the height-two witness \((k,\Gamma_2)\). At the empty parameter set we
assign the terminal presheaf value; no LKN assertion for the zero ring is
needed.

Over the perfect field \(F\), take the connected--étale sequence and its
unique splitting:
\[
 \sigma_F:\quad
 0\longrightarrow C_F\longrightarrow\Gamma_F
 \longrightarrow Q_F\longrightarrow0,
 \qquad \Gamma_F\simeq C_F\oplus Q_F .
\]
Here \(Q_F\) is étale of height one. Existence and uniqueness of the
splitting are
\cite[Proposition~2.5.20 and Remark~2.5.24]{LurieEllipticII2018}.
We apply the splitting remark only over the perfect field \(F\).
Write \(\iota_F:C_F\to\Gamma_F\) for the inclusion, and let
\[
 \lambda_0:\Gamma_F^\circ\xrightarrow{\sim}(\Gamma_2)_F^\circ
\]
be exactly the tautological formal isomorphism of the isomorphism torsor
in \cite[Example~4.6.3]{BMR2026}. Over the discrete field \(F\), the
hypotheses of
\cite[Theorem~2.3.12 and Corollary~2.3.13]{LurieEllipticII2018} hold because
\(p=0\). Their fully faithful identity-component functor gives a unique
isomorphism
\(\lambda_F:(\Gamma_2)_F\xrightarrow{\sim}C_F\), characterized by
\[
 \lambda_0\iota_F^\circ\lambda_F^\circ=\mathrm{id}.
\]
Thus our \(\lambda_F\) has the direction from \(\Gamma_2\) to the
connected summand. Define \(\sigma_T\), its splitting, \(\iota_T\), and
\(\lambda_T\) by base change to \(R_T\). In particular, \(C_T\) denotes
a \(p\)-divisible group, whereas the identity component of \(\Gamma_T\) denotes its formal group.

\paragraph{Naturality of the split reference pair.}
We check the family naturality of this splitting. If \(M\) is a finite
projective \(R_T\)-module, then
\[
 M\longrightarrow
 \prod_{t\in T}\bigl(M\otimes_{R_T,\operatorname{ev}_t}F\bigr)
\]
is injective: embed \(M\) as a direct summand of \(R_T^N\) and use
\(R_T\hookrightarrow\prod_tF\). It follows that maps between finite flat
group schemes over \(R_T\) are detected by these fibers, by applying the
same observation to the target of the corresponding coordinate-algebra
maps. It detects maps of \(p\)-divisible groups on each \(p^m\)-torsion
level. These are actual equalities of maps: finite flat algebras over the
ordinary ring \(R_T\) are ordinary finite projective algebras, and their
mapping spaces are discrete. The finite-flat functor description recalled
in the proof of \cite[Proposition~2.5.20]{LurieEllipticII2018} identifies
the compatible torsion-level maps with maps of \(p\)-divisible groups.

Let \(\rho:R_T\to R_U\) be a parameter pullback or a map furnished by a
continuous action family, and let
\(\gamma:\rho^*\Gamma_T\xrightarrow{\sim}\Gamma_U\) be the given pair
isomorphism. Both induced sequences are connected--étale: connectedness
is preserved by base change by the locally nilpotent augmentation
criterion of \cite[Remark~2.3.8]{LurieEllipticII2018}, and étaleness is
preserved by base change. The discrete characteristic-\(p\) ring \(R_U\)
is complete with ideal of definition zero. Thus
\cite[Theorem~2.5.13]{LurieEllipticII2018} gives the unique isomorphism of
these sequences extending \(\gamma\). The two resulting sections
\(Q_U\to\Gamma_U\) agree at each \(u\in U\), by uniqueness of the splitting
over the perfect field \(F\); the preceding detection argument makes them
equal over \(R_U\). Applying this with \(U=T\times G^a\), for every
\(a\geq0\), proves compatibility with all family, multiplication, and unit
maps. The same holds for \(\lambda_T\) by the fully faithful
identity-component functor. The input is the continuous action on the
full completed pair and its tautological formal trivialization from
\cite[Examples~4.6.3--4.6.4]{BMR2026}, restricted to the nonextended
\(P_3\times P_2\), which fixes \(k\). This argument uses the pulled-back
field splitting; it does not assert a splitting theorem for arbitrary
perfect rings.

\paragraph{The section on oriented deformations.}
We now construct the section on oriented deformations. Let \(A\) be any
complete adic \(E_\infty\)-ring, without a Noetherian or connectivity
assumption. For nonconnective \(A\), the notation \(\mathrm{BT}_p(A)\)
means \(\mathrm{BT}_p(\tau_{\geq0}A)\). For a finitely generated ideal of
definition \(I\), put \(B_I=\pi_0A/I\), write \(H(B_I)\) for its discrete
Eilenberg--Mac Lane \(E_\infty\)-ring, and use the canonical reduction
\[
 \tau_{\geq0}A\longrightarrow H(B_I).
\]
A subscript \(I\) below denotes this reduction or the corresponding
ordinary base change of a reference group. No map \(A\to H(B_I)\) is
asserted. Orientations remain over
the full ring \(A\): their Bott map is an equivalence of \(A\)-modules.
They are not orientations over its connective cover; see
\cite[Warning~4.3.12]{LurieEllipticII2018}.

Represent a connected oriented deformation by
\((C_A,I,\mu,\alpha_C,e)\), where \(I\subset\pi_0A\) is a finitely
generated ideal of definition,
\[
 \mu:R_T\longrightarrow\pi_0A/I,\qquad
 \alpha_C:\mu^*C_T\xrightarrow{\sim}(C_A)_{\pi_0A/I},
\]
and \(e\) orients \(C_A^\circ\) over \(A\). Tags are identified only after
passage to a common larger finitely generated ideal of definition on
which both the residue maps and group isomorphisms agree, as in
\cite[Definitions~3.1.1 and~3.1.4]{LurieEllipticII2018}.
Since \(R_T\) has characteristic \(p\), the tag forces \(p\in I\).
Completeness for \(I\) therefore implies \((p)\)-completeness, and \(p\)
is topologically nilpotent, by
\cite[Remarks~0.0.10--0.0.12]{LurieEllipticII2018}. The tag makes \(C_A\)
formally connected. These facts meet the hypotheses of the
identity-component and connected--étale results used below; they do not
assert completeness of \(\tau_{\geq0}A\).

The reduction equivalence
\[
 \mathrm{BT}^{\mathrm{et}}_p(A)\xrightarrow{\sim}
 \mathrm{BT}^{\mathrm{et}}_p(\pi_0A/I)
\]
of \cite[Corollary~2.5.10]{LurieEllipticII2018} lifts \(\mu^*Q_T\) to an
étale group \(Q_A\), with a specified residual isomorphism \(\alpha_Q\).
The space of such lifts with this isomorphism is contractible; the
unmarked object \(Q_A\) itself may have automorphisms. More precisely, if
\(\mathscr D_{C,I}^{\mathrm{or}}(A)\) denotes the fixed-\(I\) stage, use
the homotopy pullback
\[
 \mathscr X_I(A)=
 \mathscr D_{C,I}^{\mathrm{or}}(A)
 \mathop{\times}_{\mathrm{BT}^{\mathrm{et}}_p(B_I)^\simeq}
 \mathrm{BT}^{\mathrm{et}}_p(A)^\simeq ,
\]
where the first map sends a tag to \(\mu^*Q_T\) and the second is the typed
reduction above. Its projection to
\(\mathscr D_{C,I}^{\mathrm{or}}(A)\) is an equivalence by
Corollary~2.5.10; its fiber includes the marking \(\alpha_Q\).

These projections and the split construction form natural diagrams on
the category of compatible data \((A,I,\mu)\), including adic maps,
residue maps after ideal refinement, and compatible maps of split base
sequences. For an adic map \(f:A\to A'\), start with a target ideal of
definition \(J_0\). Continuity gives \(f(I^n)\subseteq J_0\) for some
\(n\), so \(J_0+(f(I))\) is again a finitely generated ideal of definition
and contains \(f(I)\). Such choices admit common larger ideals. For
\(I\subseteq J\), the same \(Q_A\) with the reduced marking is used.

The fixed-\(I\) tagged categories for complete \(A\) are groupoidal by
the proof of \cite[Lemma~3.1.10]{LurieEllipticII2018}: a tagged morphism
is an equivalence modulo \(I\), and equivalence is detected on its finite
flat \(p\)-torsion level. Orientations are pulled back from
\(\mathrm{BT}_p(A)^\simeq\), independently of \(I\); filtered colimits
of spaces commute with this pullback. Hence the filtered colimit of the
fixed stages is precisely the oriented tagging anima of
Definition~3.1.4. The auxiliary projections are an objectwise
equivalence of diagrams; take their filtered colimits and invert that
equivalence in the functor category. This supplies all lift, refinement,
and base-map coherences. We do not replace that colimit by reduction
to the nilradical, which \cite[Warning~3.1.9]{LurieEllipticII2018} excludes
in this generality.

Form the split sequence
\[
 0\longrightarrow C_A\longrightarrow C_A\oplus Q_A
 \longrightarrow Q_A\longrightarrow0.
\]
It is an exact sequence of \(p\)-divisible groups by
\cite[Proposition~2.4.8]{LurieEllipticII2018}, applied to the split pushout
square. Tag the entire residual sequence using
\[
 \mu^*\Gamma_T\simeq\mu^*C_T\oplus\mu^*Q_T
 \xrightarrow{\alpha_C\oplus\alpha_Q}
 (C_A\oplus Q_A)_{\pi_0A/I}.
\]
This retains the original full tag, including \(\mu\), and chooses no
basis or trivialization of the étale factor. Since \(A\) is
\((p)\)-complete, \(Q_A^\circ=0\) by
\cite[Proposition~2.5.8]{LurieEllipticII2018}. The identity-component fiber
sequence, as in the proof of
\cite[Proposition~2.5.17]{LurieEllipticII2018}, gives
\(C_A^\circ\xrightarrow{\sim}(C_A\oplus Q_A)^\circ\).
Transport the same orientation \(e\) through this equivalence over \(A\).

Write \(\mathscr D_\sigma^{\mathrm{or}}\),
\(\mathscr D_\Gamma^{\mathrm{or}}\), and
\(\mathscr D_C^{\mathrm{or}}\) for the resulting oriented deformation
functors. Let \(m:\mathscr D_\sigma^{\mathrm{or}}\to
\mathscr D_\Gamma^{\mathrm{or}}\) and
\(q:\mathscr D_\sigma^{\mathrm{or}}\to\mathscr D_C^{\mathrm{or}}\)
forget to the middle and left terms. Proposition~6.2.2 of
\cite{LurieEllipticII2018} makes the underlying \(m\) an equivalence.
It remains an equivalence with orientations: a tagged quotient is étale
by Remark~2.5.11 of that paper, and the left and middle identity formal
groups are naturally equivalent. Set \(\mathfrak f=qm^{-1}\), taking the
inverse in the functor category. The split construction gives
\(\widetilde{\mathfrak s}:\mathscr D_C^{\mathrm{or}}\to
\mathscr D_\sigma^{\mathrm{or}}\), whose left projection is the identity
on \((I,\mu,\alpha_C,e)\). Therefore
\[
 \mathfrak f\mathfrak s\simeq\mathrm{id},
 \qquad \mathfrak s=m\widetilde{\mathfrak s},
\]
naturally with all the preceding refinements and base maps.

\paragraph{Corepresentability and the reverse ring map.}
We justify the oriented mapping property at the possibly
non-Noetherian \(R_T\). The oriented mapping property in
\cite[Remark~6.0.7]{LurieEllipticII2018} is stated in the Noetherian
setting of its Theorem~6.0.3; the general statement of BMR
Remark~4.1.22 needs the following argument here. For either pair
\((R_T,\Gamma_T)\) or \((R_T,C_T)\), the characteristic is \(p\), and
\cite[Remark~3.4.2]{LurieEllipticII2018} supplies nonstationarity and an
almost perfect absolute cotangent complex. Theorem~3.4.1 of that paper
therefore gives a complete connective universal spectral deformation
ring \(R^{\mathrm{un}}\), a surjection
\(\pi_0R^{\mathrm{un}}\to R_T\), and a finitely generated kernel \(J\)
which is an ideal of definition.

Let \(E\) be the orientation classifier of its universal identity formal
group, as in \cite[Construction~4.1.17 and Definition~4.1.20]{BMR2026}.
The LKN witnesses and \cite[Definition~4.2.1 and Theorem~4.3.8]{BMR2026}
supply balancedness. By Remark~4.1.16 and Definitions~4.1.19--4.1.20
there, the map \(\pi_0R^{\mathrm{un}}\to\pi_0E\) is an isomorphism, the
odd homotopy groups of \(E\) vanish, and every even homotopy group is an
invertible \(\pi_0E\)-module. Here completeness of modules means
derived completeness. The homotopy-group completeness argument
in the proof of \cite[Corollary~6.0.5]{LurieEllipticII2018} now applies
without a Noetherian step. Specifically, the homotopy-group criterion
cited there as SAG Theorem~7.3.4.1 makes \(\pi_0R^{\mathrm{un}}\) \(J\)-complete because
\(R^{\mathrm{un}}\) is \(J\)-complete. An invertible module over this
ring is a direct summand of a finite free module, hence is also
\(J\)-complete. The same criterion then makes \(E\) \(J\)-complete.

Endow \(\pi_0E\) with the topology transported from
\(\pi_0R^{\mathrm{un}}\). An \(E_\infty\)-map \(E\to A\) is continuous
exactly when its restriction from \(R^{\mathrm{un}}\) is continuous,
because these \(\pi_0\) rings and topologies agree. The universal
deformation mapping property of Theorem~3.4.1, followed by the
orientation-classifier mapping property, therefore identifies the adic
mapping anima from \(E\) with the triples consisting of a deformation,
its strict tagging class, and an orientation over full \(A\). This
locally proves the use of BMR Remark~4.1.22 at \(R_T\).
It turns \(\mathfrak f\) and \(\mathfrak s\) into continuous adic maps
\[
 i_T:E(R_T,C_T)\longrightarrow E(R_T,\Gamma_T),\qquad
 r_T:E(R_T,\Gamma_T)\longrightarrow E(R_T,C_T),
 \qquad r_Ti_T\simeq\mathrm{id}.
\]
They preserve the \(SW(k)\)-structure because they retain the restriction
of the same tag \(\mu\) to \(k\). The construction is natural in split
base sequences. The family argument above therefore makes the \(r_T\)
continuous \(G\)-natural maps.

For precision, this last assertion also survives condensation. Take the
category of split LKN sequences with both full and connected pairs LKN.
The preceding construction is a natural transformation from its full
\(E\)-functor to its connected \(E\)-functor. Extend these functors and
the transformation from the presheaf category and tensor with condensed
anima as in \cite[Construction~4.5.2]{BMR2026}. Hypersheafifying the raw
split parameter diagram gives a condensed object with its continuous
\(G\)-action; its full and connected images are the original pair
diagrams. The resulting enriched transformation is the pointwise \(r_T\)
followed by hypersheafification, giving \(r^{\mathrm{cond}}\).

\paragraph{The connected tensor factor.}
It remains to identify the \(E_2\) factor. For each \(T\), let
\(b_T:E_2\to E(R_T,\Gamma_T)\) be the ordinary tensor-factor map followed
by completion, and let \(c_T:E_2\to E(R_T,C_T)\) be the canonical
connected map using \(\lambda_T\). The pointwise tensor construction uses
the Noetherian perfect reference pair \((k,\Gamma_2)\), with connected
height-two group, and the localized height-two LKN pair as its other
input. These meet \cite[Construction~4.4.7]{BMR2026}. The proof of
Theorem~4.4.8 there identifies its mapping spaces on \(K(2)\)-local
complete adic \(SW(k)\)-algebra tests. Choose a common representative
ideal \(I\) for the two \(k\)-linear tag classes, using the strict
refinement relation. If \(e_1,e_2\) orient the formal groups of the two
deformations \(H_1,H_2\), let
\[
 q_{e_i}:\Gamma_A^{Q}\xrightarrow{\sim}H_i^\circ,\qquad
 j=q_{e_2}q_{e_1}^{-1}:H_1^\circ\xrightarrow{\sim}H_2^\circ
\]
be the isomorphisms from the same Quillen formal group and their
comparison. This is the natural orientation correspondence of
\cite[Remark~4.1.18]{BMR2026}, using Lurie's Propositions~4.3.21 and
4.3.23. The orientations are over full \(A\), while \(j\) is a formal
isomorphism over \(\tau_{\geq0}A\). Only \(j\), not an orientation, is
reduced along \(\tau_{\geq0}A\to H(B_I)\).

Let \(\alpha:\Gamma_{T,I}\xrightarrow{\sim}(H_1)_I\) be the full tag
after completion, and let \(\lambda_{0,I}\) be the base change of the
same torsor isomorphism \(\lambda_0\) to \(B_I\). The tuple repacking of
Theorem~4.4.8 first goes from the reference \(\Gamma_2^\circ\) through
\(\lambda_{0,I}^{-1}\) to the reference full formal group, then through
\(\alpha^\circ\) to \((H_1)_I^\circ\), and finally through \(j_I\).
Thus the formal tag of the second tensor factor is exactly
\[
 \beta_{\mathrm{out}}^\circ
   =j_I\,\alpha^\circ\,\lambda_{0,I}^{-1}.
\]

For the split section take
\[
 H_1=C_A\oplus Q_A,\quad H_2=C_A,\quad
 e_1=(i^\circ)_*e,\quad e_2=e,
\]
where \(i:C_A\to C_A\oplus Q_A\) is the inclusion. Naturality of the
Quillen correspondence gives \(q_{e_1}=i^\circ q_e\), so
\(j=(i^\circ)^{-1}\). In our direction for \(\lambda_T\), put
\[
 \beta=\alpha_C\lambda_{T,I}:
       (\Gamma_2)_{B_I}\xrightarrow{\sim}(C_A)_I .
\]
The defining relation for \(\lambda_F\) gives
\(\lambda_{0,I}^{-1}=\iota_{T,I}^\circ\lambda_{T,I}^\circ\), and the
definition of the full split tag gives
\(\alpha^\circ\iota_{T,I}^\circ=i_I^\circ\alpha_C^\circ\).
Substitution therefore cancels the precise torsor and orientation
isomorphisms:
\[
 \begin{aligned}
 \beta_{\mathrm{out}}^\circ
  &=(i_I^\circ)^{-1}\alpha^\circ
       \iota_{T,I}^\circ\lambda_{T,I}^\circ\\
  &=(i_I^\circ)^{-1}i_I^\circ
       \alpha_C^\circ\lambda_{T,I}^\circ
    =\beta^\circ .
 \end{aligned}
\]
Theorem~2.3.12 and Corollary~2.3.13 of \cite{LurieEllipticII2018}, over
\(A\) and over the discrete characteristic-\(p\) ring \(B_I\), identify
the formally connected output with \(C_A\) and the tag itself with
\(\beta\). Their hypotheses were checked above. The output orientation
is \(e\), since pushforward through \(i^\circ\) is undone through its
inverse. The output ring tag is exactly
\(k\to R_T\xrightarrow{\mu}B_I\). Completion only precomposes \(\mu\)
and base changes the same \(\lambda_0\), so the calculation respects the
\(SW(k)\)-linear tagging condition. There is no additional Bott unit or
automorphism normalization. It is natural under strict ideal refinement
and the parameter and family maps already treated. Hence, on the stated
local test category,
\[
 (r_Tb_T)^*=b_T^*\mathfrak s_T\simeq c_T^*.
\]
The connected endpoint is \(K(2)\)-local by
\cite[Proposition~4.4.2]{BMR2026} for its height-two witness, and it is
complete by the argument above. The same statements hold for \(E_2\).
Thus both are \(K(2)\)-local complete adic \(SW(k)\)-algebras.
Yoneda on the stated test category, equivalently evaluation on the universal connected deformation at
\(A=E(R_T,C_T)\), gives \(r_Tb_T\simeq c_T\). This preserves the full
connected tag \((I,\mu,\alpha_C)\), its restricted reference-pair tag
\(\beta:(\Gamma_2)_{B_I}\xrightarrow{\sim}(C_A)_I\) along
\(k\to R_T\xrightarrow{\mu}B_I\), and the orientation \(e\); it uses
no locality of \(E(R_T,\Gamma_T)\).


\paragraph{The continuous solid identity.}
We identify the actual \(E_2\) tensor composite with \(c^\square\).
Let \(b^\square:E_2^\square\to
E^{\mathrm{cond}}(\widehat L,\Gamma^{\mathrm{un}}_{3,\widehat L})\)
be the tensor-factor map followed by completion. There are two points in
passing from the pointwise identity to solid theories. First,
the pointwise presheaf in the proof of
\cite[Proposition~4.5.12]{BMR2026} uses the constant ordinary \(E_2\), not \(E_2^\square(T)\) as the reference input
of Theorem~4.4.8. The preceding identity, natural in \(T\), therefore
identifies the two maps after precomposition from \(E_2^\delta\) and
hypersheafification. By \cite[Lemma~4.5.8]{BMR2026},
\[
 \ell:E_2^\delta\longrightarrow E_2^\square
       \simeq L^\square_{K(2)}E_2^\delta
\]
is the localization unit. The connected target is
\(L^\square_{K(2)}\)-local by its profinite evaluations, so its localization
mapping property identifies the two underlying maps from \(E_2^\square\).

Second, this does not give \(E_2^\delta\) a continuous \(P_2\)-action.
Put \(X=E_2^\square\), \(Y=\widehat B^\square\), and write \(Y^S\) for the
profinite cotensor, so \(Y^S(T)=Y(S\times T)\). For the underlying solid
spectra, the space of continuous \(G\)-equivariant maps is the bar
totalization with degree \(a\)
\[
 \operatorname{Map}(X,Y^{G^a}).
\]
Every target cotensor is local. Restriction along \(\ell\) is therefore
an equivalence of mapping spaces in each degree; transport the bar cofaces
through these equivalences. This requires no action on \(E_2^\delta\).

To identify the transported source coface, let \(g:S\to P_2\) be a
continuous family. Since \(k\) is finite, compactness gives
\(C_{\cts}(S\times T,k)=C_{\cts}(T,C_{\cts}(S,k)_{\mathrm{disc}})\).
Thus the cotensor of \(k^\delta\) by \(S\) is
\((k^\delta(S))^\delta\), where \(k^\delta(S)=C_{\cts}(S,k)\).
The family gives a pair map
\[
 \gamma_g:(k,\Gamma_2)\longrightarrow(k^\delta(S),\Gamma_2):
\]
at every torsion level its coefficients are locally constant, as required
by continuity. The naturality of Lemma~4.5.8 and the cotensor comparison
of \cite[Remark~4.5.6]{BMR2026} give the commuting square
\[
 \begin{tikzcd}
 E_2^\delta \arrow[r,"\ell"]
   \arrow[d,"E(\gamma_g)^\delta"'] &
 X \arrow[d,"a_{X,g}"]\\
 E(k^\delta(S),\Gamma_2)^\delta \arrow[r,"\ell_S"'] & X^S .
 \end{tikzcd}
\]
Here \(a_{X,g}\) is the given family action and \(\ell_S\) is the
corresponding localization unit under the cotensor identification.
Consequently the source coface after restriction is precomposition of the
connected tag by \(\gamma_g\). Target cofaces and interior cofaces are the
full-pair family maps and parameter pullbacks. The tuple identification
and \(\mathfrak f\mathfrak s\simeq\mathrm{id}\) are natural for exactly
these operations on every \(R_{G^a\times T}\). They give compatible
homotopies in all bar degrees, hence
\[
 r^{\mathrm{cond}}b^\square\simeq c^\square
\]
as continuous \(G\)-equivariant maps of underlying solid spectra.
Both sides are independently maps of condensed \(E_\infty\)-rings, and
their source and target are solid. In particular they induce the same
coefficient and even-line maps. This uses the localization mapping
property degree by degree, not an interchange of localization with
totalization or fixed points.

\end{proof}
